\title{Direct Preference Optimization: Your Language Model is Secretly a Reward Model}

\begin{abstract}
Large-scale unsupervised language models (LMs) acquire extensive knowledge about the world and develop some reasoning abilities, yet precisely guiding their outputs remains a challenge due to the unsupervised approach in their initial training. Common strategies to improve the controllability of LMs involve gathering human judgments on the relative merits of model outputs, followed by adjusting the LM via fine-tuning to better reflect these preferences—most notably through reinforcement learning from human feedback (RLHF). Nevertheless, RLHF is often intricate and unstable: it requires first training a reward model that captures human preferences, and then using reinforcement learning to adapt the LM, seeking to maximize this proxy reward while minimizing deviation from the pre-trained model. This work presents a novel reparameterization of the RLHF reward model, making it possible to directly derive the optimal policy in closed form. This insight allows us to address the conventional RLHF objective using a straightforward classification loss, resulting in an algorithm termed \textit{Direct Preference Optimization} (DPO). DPO is robust, efficient, and computationally simple, dispensing with the need for LM sampling during fine-tuning or extensive hyperparameter adjustments. Empirical results indicate that DPO fine-tunes LMs to better align with human preferences, achieving performance on par with or superior to existing approaches. Specifically, DPO outperforms RLHF methods based on PPO in sentiment control and achieves equal or higher response quality for summarization and single-turn dialogue, all while offering significant reductions in complexity and resource demands.
\end{abstract}

\section{Introduction}
Large-scale unsupervised language models (LMs), when trained on extensive corpora, demonstrate remarkable and often unexpected capabilities~\citep{chowdhery2022palm, brown2020language, touvron2023llama,bubeck2023sparks}. Despite these achievements, the underlying data originate from humans with highly diverse intentions, expertise, and priorities. This diversity means that not all learned behaviors or knowledge are desirable for AI systems to emulate; for example, although it is beneficial for an AI coding assistant to \textit{recognize} typical programming errors to enable their correction, we prefer the generated code to reflect the superior, albeit infrequent, coding proficiency present within the training data. Similarly, we expect a language model to \textit{recognize} prevalent misconceptions—such as those held by half of the population—while avoiding asserting these misconceptions as truth in half of its relevant outputs. Thus, being able to control the model’s \emph{targeted outputs and conduct} amid its broad \textit{knowledge and skillset} is fundamental for creating AI systems that are reliable, high-performing, and easy to direct \citep{ouyang2022training}. Conventional approaches typically employ reinforcement learning (RL) to align LMs with human preferences, but in this work, we demonstrate that the prevailing RL-based objective can actually be achieved precisely by adopting a straightforward binary cross-entropy loss, thereby streamlining the overall preference learning process.

Generally, prevalent techniques for preference alignment introduce preferred behaviors into language models through carefully selected datasets that encode human preferences for helpfulness and safety. This phase of preference adaptation follows large-scale, unsupervised pre-training on a broad corpus. While supervised fine-tuning on high-quality human-produced responses offers the simplest route to preference learning, methods such as reinforcement learning from human or AI feedback (RLHF/RLAIF; \citep{christiano2017deep,bai2022constitutional}) have proved the most effective. RLHF involves training a reward model to capture human preferences, followed by using RL to adapt the language model to maximize these rewards without substantial deviation from the pre-trained model. Although RLHF has achieved notable success in producing conversational and coding LMs, its learning process is significantly more intricate compared to supervised techniques: it requires the training of multiple models and repeated sampling from the language model’s policy during optimization, which substantially increases computational demands.

In this work, we demonstrate a method for training a language model to conform to human preferences without relying on explicit reward modeling or reinforcement learning approaches. We introduce \textit{Direct Preference Optimization (DPO)}, a method that, while implicitly targeting the same objective as conventional RLHF frameworks—namely, maximizing reward subject to a KL-divergence regularization—is both easier to implement and more direct to train. DPO operates by raising the log-likelihood of preferred responses relative to disfavored ones but incorporates a dynamically computed, example-specific importance weight that mitigates the model collapse observed when using a straightforward probability ratio formulation. Similar to existing algorithms, DPO makes use of a theoretical preference model (for instance, the Bradley-Terry model; \cite{bradley1952rankanalysis}) to quantify the alignment between a reward function and human preference data. Yet, unlike prior approaches that use such models to derive a preference loss for reward model training before optimizing a policy to maximize the learned reward, DPO reformulates the preference loss as a function of the policy itself via a change of variables. With a collection of annotated human preference data for model outputs, DPO enables policy optimization through a simple binary cross-entropy loss, yielding the optimal policy for an implicit reward function that captures the given preference data.

Our primary contribution is the introduction of Direct Preference Optimization (DPO), a straightforward approach for preference-driven language model training that eliminates the need for reinforcement learning. Experimental results indicate that DPO matches or exceeds the performance of existing preference learning techniques—including those based on PPO-driven RLHF—in applications such as sentiment steering, text summarization, and conversational modeling, with language models as large as 6B parameters.

\section{Related Work}

Progressively larger self-supervised language models have demonstrated the capacity to solve certain tasks in a zero-shot manner \citep{radford2019language} or with minimal prompt examples \citep{gpt3,megatron,chowdhery2022palm}. Nevertheless, their effectiveness on specific downstream applications and their ability to align with user intent are substantially enhanced through fine-tuning on instruction-based datasets paired with human-generated responses \citep{mishra-etal-2022-cross,sanh2022multitask,chung2022scaling,thoppilan2022lamda}. This process, referred to as `instruction-tuning', allows large language models (LLMs) to extend their capacity for following instructions beyond those encountered during fine-tuning, thus making them more generally useful \citep{chung2022scaling}. While instruction tuning has achieved notable advances, collecting \textit{relative} human preference data—comparing the quality of different responses—is often more feasible than obtaining expert-level demonstrations. As a result, recent research has focused on fine-tuning LLMs using datasets based on human preference comparisons, leading to improvements in translation \citep{kreutzer-etal-2018-reliability}, summarization \citep{stiennon2022learning,ziegler2020finetuning}, narrative generation \citep{ziegler2020finetuning}, and instruction adherence \citep{ouyang2022training,ramamurthy2023is}. The typical approach involves first training a neural reward model to capture the observed preferences—frequently leveraging frameworks such as the Bradley-Terry model \citep{bradley1952rankanalysis}—and subsequently using reinforcement learning techniques to adjust the language model parameters for higher reward. Common algorithms for this stage include REINFORCE \citep{williams1992reinforce}, proximal policy optimization (PPO; \cite{schulman2017proximal}), and their derivatives \citep{ramamurthy2023is}. In related research, LLMs fine-tuned with human feedback for following instructions have also been utilized to produce synthetic preference datasets targeting specific properties such as safety or harmlessness \citep{bai2022constitutional}, with only minimal human oversight provided through textual guidelines for the LLM’s evaluations. This area synthesizes advancements in two domains: one devoted to the use of reinforcement learning for optimizing language models to achieve diverse objectives~\citep{Ranzato2015SequenceLT,paulus2018a,wu2018learning}, and another focused on general methodologies for learning from human preferences \citep{christiano2017deep,kupcsik2018learning}. Despite the potential benefits of employing relative human preferences, effectively fine-tuning large language models with reinforcement learning presents significant practical hurdles; the present work offers a theoretically grounded framework for optimizing relative preferences without relying on reinforcement learning techniques.

Beyond the realm of language, the problem of deriving policies from preferences has been explored within both bandit and reinforcement learning paradigms, yielding a variety of methodologies. In contextual bandit frameworks, learning that leverages preferences or action rankings instead of explicit rewards is referred to as the contextual dueling bandit (CDB; \cite{yue2012karmed,dudik2015contextual}). In scenarios lacking absolute reward signals, theoretical investigations into CDBs replace the conventional concept of an optimal policy with the so-called \textit{von Neumann winner}, which denotes a policy that secures an expected win rate of at least 50\% against every alternative policy \citep{dudik2015contextual}. Notably, CDBs assume that preference feedback is provided interactively in an online manner, whereas in human preference learning, the standard procedure typically involves a static dataset containing preference-annotated action pairs, curated offline \citep{yan2022human}. In a similar vein, \textit{preference-based RL} (PbRL) operates using binary preferences generated according to an \textit{unknown} scoring function, as opposed to explicit rewards \citep{BusaFekete2014,ruiz2023dueling}. There exists a spectrum of PbRL algorithms, some of which are capable of leveraging off-policy preference datasets; however, these commonly entail a two-stage process involving first the explicit estimation of the underlying scoring function (the reward model), followed by policy optimization with respect to this estimate \citep{jain2013learning,BusaFekete2014,christiano2017deep,sadigh2017active,kupcsik2018learning}. By contrast, we propose a unified, single-stage approach to policy learning that directly maximizes compliance with observed preferences.

\section{Preliminaries}\label{section:prelims}

We briefly outline the RLHF framework as described in \citeauthor{ziegler2020finetuning} and subsequent works (\citep{stiennon2022learning, bai2022training, ouyang2022training}), which generally proceeds through three main steps: (1) supervised fine-tuning (SFT); (2) collection of preference data and reward model training; and (3) reinforcement learning-based policy optimization.

\textbf{SFT}: The RLHF process commences with adapting a pre-trained language model to downstream tasks (such as summarization or dialogue) by supervised learning on curated, high-quality datasets, resulting in an initial policy $\pi^\text{SFT}$.

\textbf{Reward Modelling Phase}: During the next phase, the SFT policy is conditioned on prompts $x$, generating answer pairs $(y_1, y_2)\sim \pi^\text{SFT}(y \mid x)$. Human annotators are then tasked with indicating a preferred answer, formalized as $y_w\succ y_l \mid x$, where $y_w$ is the selected response and $y_l$ the less preferred one from the pair $(y_1, y_2)$. These preferences are assumed to be produced according to some underlying, unobserved reward model $r^*(y, x)$. Several strategies exist for preference modeling, with the Bradley-Terry (BT) model \cite{bradley1952rankanalysis} being a prevalent option (though more general frameworks such as Plackett-Luce ranking models \citep{plackett1975analysis, luce2012individual} are equally applicable when rankings over more than two responses are available). The BT model represents the distribution of human preferences $p^*$ as follows:

\begin{equation}\label{eq:bradley-terry}
    p^*(y_1\succ y_2 \mid x)=\frac{\exp\left(r^*(x, y_1)\right)}{\exp\left(r^*(x, y_1)\right) + \exp\left(r^*(x, y_2)\right)}.
\end{equation}

Given a fixed dataset of pairwise comparisons, $\mathcal{D}=\bigl\{x^{(i)}, y_w^{(i)}, y_l^{(i)}\bigr\}_{i=1}^N$, which are assumed to be sampled from $p^*$, one can introduce a parameterized reward function $r_{\phi}(x, y)$ and fit its parameters by maximizing the likelihood. By interpreting this as a binary classification task, the corresponding negative log-likelihood loss takes the form:

\begin{equation}\label{eq:reward_model}
    \mathcal{L}_R(r_{\phi}, \mathcal{D}) = -\mathbb{E}_{(x, y_w, y_l)\sim \mathcal{D}}\bigl[\log \sigma(r_{\phi}(x, y_w)- r_{\phi}(x, y_l))\bigr]
\end{equation}

where $\sigma$ denotes the logistic sigmoid function. In the domain of language models, $r_{\phi}(x, y)$ is commonly initialized from the SFT model $\pi^\text{SFT}(y \mid x)$, with an additional linear head appended atop the transformer's final layer, producing a single reward scalar \cite{ziegler2020finetuning}. To mitigate variance in the reward outputs, existing literature often normalizes the rewards, enforcing that $\mathbb{E}_{x,y\sim \mathcal{D}}\left[r_\phi(x, y)\right] = 0$ across all $x$.

\textbf{RL Fine-Tuning Phase}: At the reinforcement learning stage, the acquired reward model is leveraged to give evaluative feedback to the language model. Following established practices~\citep{jaques2017sequence, jaques2020human}, the optimization is framed as

\begin{equation}\label{eq:RL}
\max_{\pi_{\theta}}  \mathbb{E}_{x\sim \mathcal{D}, y\sim \pi_{\theta}(y \mid x)}\bigl[r_{\phi}(x, y)\bigr] - \beta\mathbb{D}_{\textrm{KL}}\bigl[\pi_{\theta}(y\mid x)\mid \mid \pi_\text{ref}(y\mid x)\bigr],
\end{equation}

with $\beta$ acting as a regularization coefficient controlling divergence from a reference policy $\pi_\text{ref}$, which is typically set to the initial SFT model $\pi^\text{SFT}$. In practical implementations, the policy $\pi_\theta$ for the language model is initialized from $\pi^\text{SFT}$ as well. The KL-divergence constraint is crucial for two reasons: it keeps the policy close to the distribution where the reward model provides reliable assessments, and it guards against collapse to trivial or repetitive outputs, thereby supporting output diversity. As language generation is inherently discrete, this objective function cannot be differentiated directly, necessitating the use of reinforcement learning methods. The standard technique \citep{ziegler2020finetuning, stiennon2022learning, bai2022training, ouyang2022training} is to define the reward as ${r(x, y) = r_{\phi}(x, y) -\beta (\log \pi_{\theta}(y\mid x) - \log \pi_\text{ref}(y\mid x))}$ and maximize it using PPO \cite{schulman2017proximal}.

\section{Direct Preference Optimization}\label{sec:DPO}

In light of the substantial computational burden posed by reinforcement learning when scaling up to large language models, we propose to develop a streamlined method for directly optimizing policies from preference data. Departing from the conventional RLHF approach—which separates reward modeling from policy optimization via RL—our method employs a particular reward model parameterization that facilitates analytic derivation of the optimal policy, thereby removing the need for an RL-based training process. As we elaborate below, our central contribution is to exploit an explicit, closed-form correspondence between reward functions and their optimal policies. This mapping permits reformulation of a reward-based loss into a loss directly defined over policy parameters, effectively implementing a change of variables. Through this construction, we sidestep the need to explicitly fit a standalone reward function, yet our approach remains compatible with

\textbf{Derivation of the DPO objective.} We consider the same reinforcement learning objective as in previous studies, Eq.~\ref{eq:RL}, and assume a general reward function $r$. Building on established results~\citep{peters2007reinforcement, peng2019advantage, korbak2022reinforcement, go2023aligning}, the solution that optimizes the KL-regularized reward maximization (Eq.~\ref{eq:RL}) can be written as:

\begin{equation}\label{eq:op_policy}
    \pi_r(y\mid x) = \frac{1}{Z(x)}\pi_\text{ref}(y\mid x)\exp\left(\frac{1}{\beta}r(x, y)\right),
\end{equation}

where the normalization constant $Z(x) =\sum_{y}\pi_\text{ref}(y\mid x)\exp\left(\frac{1}{\beta}r(x, y)\right)$ ensures proper normalization. For the detailed derivation, see Appendix \ref{app:derivation1}. Although one might substitute an MLE-based estimate $r_\phi$ for the unknown reward $r^*$, computing $Z(x)$ remains computationally demanding~\citep{korbak2022reinforcement, go2023aligning}, limiting practical use of this policy form. Nonetheless, it is possible to reformulate Eq.~\ref{eq:op_policy} so that the reward is expressed via the optimal policy $\pi_r$, the reference $\pi_\text{ref}$, and the partition function $Z(\cdot)$. By taking the logarithm of both sides and rearranging, we obtain:

\begin{equation}\label{eq:main_eq}
    r(x,y) =\beta \log \frac{\pi_r(y\mid x)}{\pi_\text{ref}(y\mid x)} + \beta \log Z(x).
\end{equation}

Applying this transformation to the true reward $r^*$ and the associated optimal policy $\pi^*$, we note that the Bradley-Terry model depends only on differences in rewards between completions: ${p^*(y_1 \succ y_2 \mid x) = \sigma(r^*(x, y_1) - r^*(x, y_2))}$. When we substitute the representation for $r^*(x,y)$ from Eq.~\ref{eq:main_eq} into the preference model Eq.~\ref{eq:bradley-terry}, the partition function $Z(x)$ cancels out. This allows the human preference probability to be written solely as a function of $\pi^*$ and $\pi_\text{ref}$. Thus, under the Bradley-Terry assumption, the optimal RLHF policy $\pi^*$ adheres to:

\begin{equation}\label{eq:objective}
    p^*(y_1\succ y_2 \mid x)=\frac{1}{1 + \exp\left(\beta \log \frac{\pi^*(y_2\mid x)}{\pi_\text{ref}(y_2\mid x)} - \beta \log \frac{\pi^*(y_1\mid x)}{\pi_\text{ref}(y_1\mid x)}\right)}
\end{equation}

The full derivation can be found in Appendix~\ref{app:derivation2}. Although Eq.~\ref{eq:objective} is based on the Bradley-Terry framework, analogous results follow for more general preference models like the Plackett-Luce~\citep{plackett1975analysis, luce2012individual}—see Appendix~\ref{app:plackett_luce_models}.

Having now characterized the human preference likelihood in terms of the optimal policy rather than the reward, we can now propose a maximum likelihood learning objective for a parameterized policy $\pi_\theta$. In close analogy to reward-based modeling (see Eq.~\ref{eq:reward_model}), the corresponding policy learning objective is given by:

\begin{equation}\label{eq:optimum_model}
    \mathcal{L}_\text{DPO}(\pi_{\theta}; \pi_\text{ref}) = -\mathbb{E}_{(x, y_w, y_l)\sim \mathcal{D}}\left[\log \sigma \left(\beta \log \frac{\pi_{\theta}(y_w\mid x)}{\pi_\text{ref}(y_w\mid x)} - \beta \log \frac{\pi_{\theta}(y_l\mid x)}{\pi_\text{ref}(y_l\mid x)}\right)\right].
\end{equation}

In this framework, we infer an implicit reward function via a reparameterized approach, where the optimal policy is given directly by $\pi_\theta$. Additionally, because our method mirrors fitting a reparameterized Bradley-Terry model, it inherits several theoretical guarantees—such as consistency, provided certain assumptions about the distribution of preference data hold \cite{bong2022generalized}. A more detailed examination of DPO’s theoretical aspects, especially in relation to existing literature, appears in Section~\ref{sec:theory}.

\textbf{How does the DPO update operate?} To elucidate the inner workings of DPO, one can inspect the gradient of the loss, $\mathcal{L}_\text{DPO}$. The parameter gradient $\nabla_\theta \mathcal{L}_\text{DPO}$ can be expressed as:

\begin{multline*}\label{eq:gradient}
    \nabla_\theta \mathcal{L}_\text{DPO}(\pi_\theta;\pi_\text{ref}) = \\ -\beta\mathbb{E}_{(x, y_w, y_l) \sim \mathcal{D}} \bigg[\underbrace{\sigma(\hat{r}_\theta(x, y_l) - \hat{r}_\theta (x, y_w))}_\text{greater emphasis when reward estimate errs}\bigg[\underbrace{\nabla_\theta\log \pi(y_w \mid x)}_\text{promote $y_w$} - \underbrace{\nabla_\theta\log\pi(y_l \mid x)}_\text{penalize $y_l$}\bigg]\bigg],
\end{multline*}

where the implicit reward function is defined as $\hat{r}_\theta(x, y) = \beta \log \frac{\pi_\theta(y \mid x)}{\pi_\text{ref}(y \mid x)}$, involving both the learned policy $\pi_\theta$ and a reference $\pi_\text{ref}$ (see also Section~\ref{sec:theory}). Conceptually, the gradient of $\mathcal{L}_\text{DPO}$ acts to raise the probability assigned to the favored responses $y_w$ while suppressing that of the less preferred $y_l$. Crucially, each example’s contribution is weighted by how much the implicit reward model, $\hat{r}_\theta$, overrates the less preferred choice, reflecting the extent to which the model’s current ranking of completions is incorrect (modulated by $\beta$ to reflect the KL regularization). Empirically, this form of weighting proves vital; omitting the weighting factor leads to degenerate behavior in the language model (see Appendix Table~\ref{tab:unlikelihood_generations}).

\textbf{DPO overview.} The typical workflow for DPO proceeds as follows: 1) For each prompt $x$, generate completions $y_1, y_2 \sim \pi_\text{ref}(\cdot \mid x)$, assign human preference labels, and thereby build an offline preference dataset $\mathcal{D} = \{x^{(i)}, y_w^{(i)}, y_l^{(i)}\}_{i=1}^N$; 2) Adjust the language model $\pi_\theta$ to minimize $\mathcal{L}_\text{DPO}$, using the specified $\pi_\text{ref}$, dataset $\mathcal{D}$, and chosen $\beta$. 
In applied settings, rather than collecting new human annotations and samples, it is often preferable to leverage existing publicly available preference datasets. Since such datasets are typically obtained using $\pi^\text{SFT}$, we set $\pi_\text{ref} = \pi^\text{SFT}$ when this is feasible. If $\pi^\text{SFT}$ is not accessible, $\pi_\text{ref}$ is instead initialized by maximizing the likelihood of the preferred completions $(x, y_w)$; specifically, we choose ${\pi_\text{ref} = \argmax_{\pi}\mathbb{E}_{x, y_w \sim \mathcal{D}}\left[\log \pi(y_w \mid x)\right]}$. This approach helps address potential distribution mismatch between the ideal (but unknown) reference distribution and the operational $\pi_\text{ref}$ within DPO. Further implementation details and information on hyperparameter choices are discussed in Appendix~\ref{app:implementation}.

\section{Theoretical Analysis of DPO} In this section, we provide a deeper conceptual understanding of DPO, offer theoretical justification, and highlight benefits of DPO in relation to challenges found in actor-critic methods for RLHF (for example, PPO~\cite{schulman2017proximal}).

\label{sec:theory}

\subsection{Your Language Model Is Secretly a Reward Model} DPO achieves the combined goals of reward learning and policy improvement using a unified maximum likelihood objective, thereby sidestepping both explicit reward fitting and reinforcement learning stages. Importantly, the objective in Eq. \ref{eq:main_eq} matches the Bradley-Terry model structure, where the reward is parameterized as $r^*(x, y) = \beta \log\frac{\pi^*_\theta(y \mid x)}{\pi_\text{ref}(y \mid x)}$. The optimization is then performed over $\pi_\theta$, analogous to the reward model optimization from Eq. \ref{eq:reward_model}, given an appropriate variable change. This section develops the theoretical framework for this reparameterization, demonstrates that it maintains full flexibility in the space of reward models, and proves that the optimal policy can be exactly recovered. We start by formalizing an equivalence relationship among reward functions.

\begin{definition}
Two reward functions $r(x, y)$ and $r'(x, y)$ are said to be equivalent if and only if there exists a function $f$ such that ${r(x, y)-r'(x, y) = f(x)}$.
\end{definition}

It is straightforward to verify that this notion indeed defines an equivalence relation, leading to a partitioning of all reward functions into distinct equivalence classes. The following two lemmas capture key properties of these classes:

\begin{lemma}\label{lemma:same_prefrence} 
Within the Plackett-Luce preference framework, including the Bradley-Terry model as a special case, any two reward functions belonging to the same equivalence class induce identical preference distributions.
\end{lemma}

\begin{lemma}\label{lemma:same_policy}
    Any pair of reward functions from the same equivalence class will yield the same optimal policy when considering the constrained reinforcement learning problem.
\end{lemma}

As the arguments are elementary, we postpone detailed proofs to Appendix \ref{app:lemma1}. The first lemma highlights a well-recognized identifiability issue that arises in the Plackett-Luce family of models \cite{plackett1975analysis}. This ambiguity often necessitates the imposition of further identifiability constraints to ensure meaningful maximum likelihood estimation, as discussed in Eq. \ref{eq:reward_model} and in \cite{bong2022generalized}. The second lemma asserts that, for the purpose of policy optimization, it suffices to recover any representative from the relevant equivalence class of rewards, as they all result in the same optimal policy. The following theorem, whose proof can be found in Appendix~\ref{app:thm1}, formalizes the structure of such classes:

\begin{theorem}\label{thm:main}
    Assuming mild technical conditions, every reward equivalence class consistent with the Plackett-Luce (and specifically Bradley-Terry) models can be represented using the reparameterization ${r(x, y) = \beta \log \frac{\pi(y\mid x)}{\pi_\text{ref}(y\mid x)}}$ for some model $\pi(y\mid x)$ and some fixed reference model $\pi_\text{ref}(y \mid x)$.
\end{theorem}

\begin{sproof}
    Let $r(x, y)$ denote any given reward function, and let $\pi_r(y \mid x)$ be its associated optimal model, as defined by Eq. \ref{eq:op_policy}. We demonstrate that one can construct an equivalent reward function from the class of $r$ that takes the reparameterized form above. To that end, define the projection $f$ as follows:
\begin{equation}
    f(r; \pi_\text{ref}, \beta)(x, y) = r(x, y) - \beta\log\sum_{y}\pi_\text{ref}(y\mid x)\exp\left(\frac{1}{\beta}r(x, y)\right)
\end{equation}
This operation normalizes the original reward by subtracting the log-partition function of $\pi_r$ with respect to $\pi_\text{ref}$. Since the normalization term depends solely on the prefix $x$, $f(r; \pi_\text{ref}, \beta)(x, y)$ remains in the same equivalence class as $r(x, y)$. By substituting $r$ with the right-hand side of Eq.~\ref{eq:main_eq} (valid for any $r$), it follows that $f(r; \pi_\text{ref}, \beta)(x, y) = \beta \log \frac{\pi_r(y\mid x)}{\pi_\text{ref}(y\mid x)}$. Thus, the projection $f$ produces an equivalently classed reward function of the desired parametric form, implying that our reward model incurs no generality loss by adopting this reparameterization.
\end{sproof}

An alternative perspective on Theorem~\ref{thm:main} is that it identifies the specific reward function, within each equivalence class, that is selected by the DPO reparameterization. More precisely, it singles out the reward function for which

\begin{equation}\label{eq:lag_p}
     \sum_{y}\underbrace{\pi_\text{ref}(y\mid x)\exp\left(\frac{1}{\beta}r(x, y)\right)}_{=\pi(y\mid x)\text{, using Thm.~\ref{thm:main} reparam.}} = 1,
\end{equation}

meaning that $\pi(y\mid x)$ constitutes a proper probability distribution, with non-negative values that sum to one. Furthermore, from Eq.~\ref{eq:op_policy}, it is evident that Eq.~\ref{eq:lag_p} corresponds to the partition function of the optimal policy associated with the reward $r(x, y)$. The principal innovation underlying the DPO approach lies in the imposition of specific constraints on the otherwise underdetermined Plackett-Luce (and, by extension, Bradley-Terry) class of preference models. This constraint preserves the expressiveness of the reward models, while ensuring that the optimal policy described by Eq.~\ref{eq:op_policy} remains analytically manageable for every prompt $x$.

\subsection{Instability of Actor-Critic Algorithms} Our framework also facilitates the identification of sources of instability in conventional actor-critic algorithms, such as PPO, widely adopted in RLHF. Focusing on the RL fine-tuning phase as detailed in Section \ref{section:prelims}, we can relate these issues to the control-as-inference paradigm \cite{levine2018reinforcement} applied to the constrained RL task presented in \ref{eq:RL}. Taking a parameterized policy $\pi_{\theta}(y\mid x)$, the goal is to minimize $\mathbb{D}_{\text{KL}}[\pi_{\theta}(y|x) \mid \mid \pi^*(y\mid x)]$, where $\pi^*$ denotes the optimal policy from Eq.~\ref{eq:optimum_model} as determined by the reward function $r_{\phi}(y, x)$. Through straightforward algebraic manipulation, this yields the objective:

\begin{equation}\label{eq:AC}
    \max_{\pi_{\theta}}\mathbb{E}_{\pi_{\theta}(y\mid x)}\bigg[\underbrace{r_{\phi}(x, y) -\beta\log\sum_{y}\pi_\text{ref}(y\mid x)\exp\left(\frac{1}{\beta}r_{\phi}(x, y)\right)}_{f(r_{\phi}, \pi_\text{ref}, \beta)} - \underbrace{\beta\log\frac{\pi_{\theta}(y\mid x)}{\pi_\text{ref}(y\mid x)}}_{\text{KL}}\bigg]
\end{equation}

This objective aligns with those targeted in earlier studies 
\citep{ziegler2020finetuning, stiennon2022learning, bai2022training, ouyang2022training}, which optimize the DPO-equivalent reward for the reward class $r_{\phi}$. In this framework, the normalization component in $f(r_{\phi}, \pi_\text{ref}, \beta)$ can be understood as the soft value function corresponding to the reference policy $\pi_\text{ref}$. Although this normalization does not influence the location of the optimum, omitting it can lead to elevated variance in the policy gradient, thereby destabilizing training. This normalization can, in principle, be handled through a learned value function; however, optimizing such a function is often challenging in practice. Alternatively, previous approaches have implemented reward normalization via a human completion baseline, effectively employing a Monte Carlo estimate with a single sample for the normalization. In contrast, the DPO reparameterization produces a reward function that obviates the need for explicit baselines. 

\section{Experiments}
Here, we empirically investigate the effectiveness of DPO for training policies directly from preference data. Our initial experiments are conducted in a controlled text-generation environment, where we compare DPO's efficiency in balancing reward maximization and KL-divergence minimization with the reference policy against standard preference learning approaches like PPO. Subsequently, we examine DPO's capabilities on more complex RLHF tasks and with larger model architectures, focusing on applications such as summarization and dialogue. Notably, we observe that DPO, with minimal hyperparameter tuning, often matches or surpasses robust baselines—including RLHF with PPO and the approach of selecting the top trajectory from $N$ samples according to a learned reward model. Before discussing these empirical findings, we outline our experimental methodology; see Appendix~\ref{app:exp_details} for further specifics.

\textbf{Tasks.} Our study investigates three distinct open-domain text generation tasks. Throughout all experiments, the methods learn a policy based on a preference dataset $\mathcal{D}=\bigl\{x^{(i)}, y_w^{(i)}, y_l^{(i)}\bigr\}_{i=1}^N$. For \textbf{controlled sentiment generation}, $x$ denotes a leading segment from an IMDb movie review \cite{maas-EtAl:2011:ACL-HLT2011}, and the aim is for the policy to generate $y$ that exhibits positive sentiment. To ensure systematic evaluation, we construct preference pairs for this task using a pre-trained sentiment classifier such that $p(\text{positive}\mid x,y_w)>p(\text{positive}\mid x,y_l)$. The SFT baseline involves fine-tuning GPT-2-large on the IMDb training split until convergence (details in App~\ref{app:sentiment_details}). For \textbf{summarization}, $x$ is a Reddit forum thread, and the task requires generating a summary $y$ capturing key information. In line with established protocols, we use the Reddit TL;DR summarization dataset \citep{volske-etal-2017-tl} together with human preference data collected by \citeauthor{stiennon2022learning}. The SFT model is fine-tuned on summaries written by humans\footnote{\url{https://huggingface.co/CarperAI/openai_summarize_tldr_sft}} within the TRLX \citep{leandro_von_werra_2023_7790115} RLHF framework. Here, the human preference set originates from \citeauthor{stiennon2022learning}, who evaluated outputs from a comparable SFT model. For \textbf{single-turn dialogue}, $x$ represents a human prompt, which could range from scientific questions to personal requests. The policy’s objective is to provide a relevant and constructive reply $y$. We utilize the Anthropic Helpful and Harmless dialogue dataset \citep{bai2022training}, comprising 170k human-assistant conversations. Each record concludes with two model-generated responses and an annotation marking the human-preferred option. Since a dedicated SFT model is unavailable in this case, we instead fine-tune a general-purpose language model using only the preferred completions to derive the SFT baseline.

\textbf{Evaluation.} We employ two distinct strategies for evaluation in our experiments. To rigorously assess how well each algorithm achieves the constrained reward maximization goal, we first focus on the controlled sentiment generation scenario. Here, we compute and plot the Pareto frontier of each method with respect to reward and KL-divergence from the reference policy, leveraging access to an explicit ground-truth reward signal (i.e., a sentiment classifier). Conversely, when a precise reward function is unavailable—as is often the case in practical applications—we instead measure each algorithm's \textit{win rate} compared to a baseline policy. This is done by soliciting judgments from GPT-4, which acts as a surrogate for human evaluation in both summarization and single-turn dialogue tasks; summary quality and response helpfulness are assessed accordingly. For the summarization experiments, the baseline comprises reference summaries from the test set, whereas for dialogue, it is the response identified as preferable in the test data. While prior research indicates that large language models serve as more reliable automated evaluators than many existing metrics \citep{Chen2023ExploringTU}, we additionally perform a human evaluation study to substantiate our reliance on GPT-4 for assessment purposes, as described in Sec.~\ref{sec:human-judgments}. Our findings reveal a strong alignment between GPT-4 and human raters, with the degree of human-GPT-4 agreement typically matching or surpassing inter-annotator agreement rates among humans.

\textbf{Methods.} Beyond DPO, we benchmark several established techniques for training language models to reflect human preferences. As straightforward baselines, we utilize zero-shot prompting of \textbf{GPT-J} \citep{gpt-j} for summarization and 2-shot prompting with \textbf{Pythia-2.8B} \citep{biderman2023pythia} in dialogue. Further, we include both the \textbf{SFT} model and the \textbf{Preferred-FT} model, the latter being fine-tuned via supervised learning on preferred completions $y_w$, sourced either from SFT outputs (for sentiment and summarization) or a general-purpose language model (for dialogue). Additionally, we examine the \textbf{Unlikelihood}~\citep{welleck2019neural} approach, which encourages the policy to allocate higher probability to $y_w$ and to \textit{downweight} $y_l$, with an adjustable unlikelihood loss coefficient $\alpha\in[0,1]$. We also assess \textbf{PPO} \citep{schulman2017proximal} utilizing a reward model trained on preference data, and \textbf{PPO-GT}, an oracle variant that accesses the true reward function (permissible only in controlled sentiment experiments). In the sentiment generation context, we test two PPO-GT implementations: an out-of-the-box version \cite{leandro_von_werra_2023_7790115} and a modified variant that normalizes rewards and tunes hyperparameters for enhanced results (the same modifications are used when deploying PPO with learned rewards). Lastly, we implement the \textbf{Best of $N$} baseline, which draws $N$ completions from the SFT (or Preferred-FT, for dialogue) model and selects the top response based on a reward model trained on preference annotations. While this method effectively separates reward model quality from PPO optimization, it is computationally burdensome, as inference requires generating $N$ responses per input, making it impractical for even moderate $N$.

\subsection{How well can DPO optimize the RLHF objective?}

The reward maximization objective in standard RLHF methods, which incorporates a KL constraint, is designed to maximize reward without allowing the policy to diverge substantially from the reference policy. Thus, a meaningful comparison between algorithms must jointly consider both the reward obtained and the KL divergence; an increase in reward accompanied by a disproportionate increase in KL is not generally preferable. In Figure~\ref{fig:frontier-tldr-main}, we present the trade-off frontier between reward and KL divergence across various algorithms on the sentiment task. For each algorithm, we conduct several runs, varying the policy conservativeness parameter in each instance (target KL values $\in\{3,6,9,12\}$ for PPO, $\beta \in \{0.05,0.1,1,5\}$ and $\alpha\in\{0.05,0.1,0.5,1\}$ for unlikelihood, and using different random seeds for preferred-FT), totaling 22 training runs. Policies are assessed after every 100 training steps up to convergence, evaluating the average reward under the actual reward function and the mean sequence-level KL\footnote{Specifically, we sum the KL-divergence at each timestep.} relative to the reference policy $\text{KL}\left(\pi\mid \mid \pi_\text{ref}\right)$. Our findings demonstrate that DPO produces a substantially more favorable frontier than other approaches, securing the highest reward while maintaining low KL. This outcome stands out for two main reasons. First, despite both DPO and PPO targeting the same optimization objective, DPO is significantly more sample-efficient, with its reward/KL trade-off strictly outperforming that of PPO. Second, DPO manages to establish a superior frontier to PPO even in scenarios where PPO has access to the true reward signal (PPO-GT).

\subsection{Can DPO scale to real preference datasets?}
\label{sec:dpo-real-datasets}
We proceed to assess how well DPO performs when fine-tuned on tasks involving summarization and single-turn dialogue. In summarization, commonly used automated evaluation metrics like ROUGE often show weak alignment with human judgment~\citep{stiennon2022learning}. Previous studies have demonstrated that language models fine-tuned with PPO on human preference data tend to generate more satisfactory summaries. To compare different approaches, we generate outputs from the test portion of the TL;DR summarization dataset and calculate the mean win rate relative to reference summaries within this test set. For every method, completions are sampled across a temperature range from 0.0 to 1.0, with win rates presented in Figure~\ref{fig:frontier-tldr-main} (right). DPO, PPO, and Preferred-FT all perform additional fine-tuning on the same GPT-J SFT base model\footnote{\url{https://huggingface.co/CarperAI/openai_summarize_tldr_sft}}. Our findings indicate that DPO achieves a win rate close to 61\% at a temperature of 0.0, surpassing PPO, which peaks at roughly 57\% under the same conditions. Notably, DPO also attains a higher best-case win rate than the $N$ baseline. It is important to mention that we did not conduct extensive tuning of the $\beta$ hyperparameter for DPO, suggesting that the observed results may not reflect DPO's full capabilities. Additionally, DPO exhibits greater resilience to changes in sampling temperature, whereas PPO's effectiveness drops to that of the original GPT-J at higher temperatures. Preferred-FT yields minimal improvement over the SFT baseline. Section~\ref{sec:human-judgments} presents a direct human comparison between DPO and PPO, revealing that samples from DPO at temperature 0.25 were favored 58\% of the time over PPO outputs at temperature 0.

For single-turn dialogue evaluation, we assess the performance of various approaches on the portion of the Anthropic HH test set \citep{bai2022training} that consists of a single round of human-assistant exchange. In the case of GPT-4, the reference completions are the preferred ones from the test set, which are used to calculate the win rate for each method. Since a standard SFT model does not exist for this task, our methodology involves initializing with a pre-trained Pythia-2.8B, applying Preferred-FT to adapt a reference model to the distribution of selected completions, and subsequently training using DPO. Our analysis also includes a comparison with the top choice among 128 Preferred-FT completions (the Best of $N$ baseline was found to plateau at $N=128$; see Appendix Figure~\ref{fig:best-of-n}) as well as a 2-shot prompt configuration of the Pythia-2.8B base model. The results indicate that, across the optimal temperature settings for each technique, DPO matches or surpasses their best win rates. We further assess an RLHF model using PPO, trained on the Anthropic HH dataset \footnote{\url{https://huggingface.co/reciprocate/ppo_hh_pythia-6B}} from a recognized repository \footnote{\url{https://github.com/CarperAI/trlx/tree/main/examples/hh}}, yet could not identify a prompting or sampling temperature that outperforms the raw Pythia-2.8B model. Given our findings from TL;DR and that both approaches target the same reward function, we treat Best of 128 as an approximate indicator of PPO-level capability. In summary, DPO stands out as the only method that is computationally efficient while consistently exceeding the preferred completions from the Anthropic HH data, and it achieves performance on par with or better than the much more resource-intensive Best of 128 method. As depicted in Figure~\ref{fig:dialogue-main}, DPO also attains its optimal outcomes in a relatively short number of training iterations.

\subsection{Generalization to a new input distribution}

To more rigorously assess how PPO and DPO perform when faced with shifts in input distribution, we apply the trained policies from our Reddit TL;DR summarization setup to a distinctly different data domain: news articles from the CNN/DailyMail test set \citep{nallapati-etal-2016-abstractive}. We employ the same optimal sampling temperatures identified from the TL;DR experiments (0 and 0.25) in this evaluation. Table~\ref{tab:ood} displays the outcomes. To determine model performance, we measure GPT-4 win rates against reference summaries, utilizing the identical GPT-4 (C) evaluation prompt as for Reddit TL;DR—modified only by substituting ``forum post'' with ``news article''. On this out-of-distribution task, DPO significantly surpasses PPO in effectiveness. These findings offer preliminary support that DPO can generalize at least as robustly as PPO, even though DPO does not leverage the unlabeled Reddit TL;DR prompts incorporated during PPO training.

\subsection{Validating GPT-4 judgments with human judgments}
\label{sec:human-judgments}
We run a human evaluation to check the consistency of GPT-4's scoring, drawing from outputs in the TL;DR summarization experiments and using two variations of GPT-4 prompts. The \textbf{GPT-4 (S)} (simple) prompt requests a judgment on which summary better captures the critical details of the post, while the \textbf{GPT-4 (C)} (concise) prompt adds an explicit preference for conciseness—motivated by observations that with the \textbf{GPT-4 (S)} prompt, GPT-4 often favors longer and more redundant summaries compared to human raters. The full texts of these prompts are available in Appendix~\ref{app:prompts}. Our comparative analysis spans three methods: the strongest (DPO, temp. 0.25), the weakest (PPO, temp. 1.0), and a representative middle performer (SFT, temp. 0.25), with each compared to greedily-sampled PPO at its best-performing temperature. In both prompt settings, GPT-4's preferences align with human assessments roughly as often as human raters concur among themselves, indicating that GPT-4 serves as a reliable surrogate for human evaluation (due to limited human resources, multiple human ratings were collected only for DPO and PPO-1 matchups). In summary, the \textbf{GPT-4 (C)} prompt yields win rates that more closely mirror human judgment; thus, we employ this prompt for principal analyses in Section~\ref{sec:dpo-real-datasets}. Further methodological details, including the rater web interface and the roster of participating human volunteers, are provided in Appendix~\ref{app:human-study}.

\section{Discussion}
Preference-based learning serves as an effective and scalable approach for aligning and enhancing the capabilities of language models. In this work, we have presented DPO, a streamlined training method that enables language models to learn from preferences without the need for reinforcement learning. DPO circumvents the necessity to force the preference learning challenge into a conventional RL framework for compatibility with typical RL methods. Instead, DPO establishes a correspondence between language model policies and reward functions, which allows models to be trained to reflect human preferences directly using a straightforward cross-entropy objective—eschewing both reinforcement learning and any sacrifice in generality. Remarkably, DPO achieves performance on par with or exceeding prevalent RLHF techniques such as those leveraging PPO, all while requiring minimal hyperparameter tuning. This substantially lowers the obstacles to training language models using human preferences.

\textbf{Limitations \& Future Work.} The findings presented prompt a variety of pertinent avenues for continued research. One open question concerns the generalization ability of the DPO policy on data outside the training distribution, especially when contrasted with approaches that rely on explicit reward models. Our preliminary experiments indicate that the generalization behavior of DPO-trained models resembles that of those trained via PPO, yet a more thorough investigation is warranted. For instance, it remains to be determined whether self-labeling using the DPO policy can capitalize on unlabeled prompts as effectively. Additionally, further study is needed to understand how issues like reward over-optimization arise within the direct preference optimization paradigm, and whether the modest reduction in performance observed in Figure~\ref{fig:dialogue-main}-right can be attributed to this effect. Furthermore, although our current evaluations include models with up to 6B parameters, scaling DPO to even larger, state-of-the-art architectures represents a promising future direction. In terms of evaluation methodology, our analyses reveal that GPT-4-based win rates are sensitive to the prompting procedure; thus, optimizing prompt design for reliable automated judgments is an important topic for subsequent work. Beyond language modeling, DPO holds potential for broader application, such as in the training of generative models across various modalities.

\end{document}